\documentclass[11pt,letterpaper]{article}

\usepackage{hyperref}
\usepackage{authblk}
\usepackage{amsmath,amsfonts,amssymb,graphicx}
\usepackage{algorithm}
\usepackage{algpseudocode}
\usepackage{mathrsfs}
\usepackage{amsthm}
\usepackage{bm}
\usepackage{titlesec}
\usepackage{setspace}
\usepackage{enumerate}
\usepackage{threeparttable}
\usepackage[square,comma,sort&compress,numbers]{natbib}
\usepackage{color}
\usepackage{tabularx,booktabs}
\usepackage{multirow}
\usepackage{colortbl}
\usepackage{xcolor}
\usepackage{geometry}
\usepackage{caption}
\usepackage{float}
\usepackage{url}
\usepackage{fullpage}
\usepackage{lipsum}

\newcommand{\junk}[1]{}
\geometry{left=1in,right=1in,top=1in,bottom=1in}

\newcommand{\red}[1]{\textcolor{red}{#1}}
\providecommand{\rev}[1]{\red{#1}}
\renewcommand{\thesubsection}{\arabic{subsection}}

\begin{document}
\thispagestyle{empty}

\title{{\huge \textbf{Supplementary Information}}\\
mBatchNet: an interactive web server for diagnosis, correction and benchmarking of batch effects in microbiome data}

\author[1,2]{Chentong Sun}
\author[1]{Shiyuan Wang}
\author[1]{Qiwei Zhang}
\author[3]{Ruishan Liu}
\author[1]{Yuxuan Du\thanks{Corresponding author. Email: yuxuan.du@utsa.edu}}

\affil[1]{Department of Electrical Engineering, The University of Texas at San Antonio, San Antonio, Texas, 78249, USA}
\affil[2]{College of Engineering and Applied Science, University of Colorado Boulder, Boulder, Colorado, 80309, USA}
\affil[3]{Department of Computer Science, University of Southern California, Los Angeles, California, 90089, USA}

\renewcommand*{\Affilfont}{\small}
\renewcommand\Authands{ and }

\date{}
\maketitle

\section*{Table of contents}

\noindent \textbf{Supplementary Notes}
\hfill \pageref{sec:supp-notes}\\

\noindent \textbf{Supplementary Figures}
\hfill \pageref{sec:supp-figures}\\

\noindent \textbf{Supplementary Tables}
\hfill \pageref{sec:supp-tables}\\

\noindent \textbf{Supplementary References}
\hfill \pageref{sec:supp-refs}\\


\newpage
\phantomsection
\label{sec:supp-notes}


\newpage
\subsection*{Supplementary Note 1. Analysis methods}
\label{sec:supp-note2}

The correction-method categories described in the main text are used as practical implementation labels rather than as a ranking system. By original application domain, microbiome-oriented methods include ConQuR, MMUPHin, PLSDA-batch, DEBIAS-M and MetaDICT, whereas general-purpose or expression-derived adjustment frameworks include ComBat, limma, ComBat-seq, FAbatch, BMC, FSQN and RUV-III-NB. By input scale and modeling family, some methods can belong to more than one category: count-aware or zero-inflation-aware count models include ConQuR, ComBat-seq and RUV-III-NB; empirical-Bayes, linear-model, model-based or quantile-normalization approaches include ComBat, limma, BMC, FAbatch and FSQN; target-aware or latent-component approaches include PLSDA-batch and FAbatch, and DEBIAS-M also uses phenotype-prediction parameters in its model; microbiome-focused processing-bias, meta-analysis or dictionary-learning approaches include DEBIAS-M, MMUPHin and MetaDICT.

The diagnostic categories are mapped as follows. Ordination diagnostics include PCA, PCoA and NMDS. Distance-matrix tests and summaries include ANOSIM, PERMANOVA and dissimilarity heatmaps. Variance-partitioning analyses include feature-wise ANOVA $R^2$, pRDA and PVCA. Neighborhood and embedding-based summaries include alignment score, entropy of batch mixing and target-label silhouette.

\begin{enumerate}
\renewcommand{\labelenumi}{(\roman{enumi})}

\item \textbf{Ordination diagnostics.}
We visualize dominant structure using principal component analysis (PCA) in Aitchison space and principal coordinates analysis (PCoA) and non-metric multidimensional scaling (NMDS) on Aitchison and Bray--Curtis dissimilarities. Plots are colored by batch and by target label to inspect batch mixing and phenotype separation. For NMDS, we report the stress value, a goodness-of-fit measure for low-dimensional representation, where lower stress indicates a more faithful embedding.

\item \textbf{Distance-matrix tests and dissimilarity summaries.}
We quantify residual batch association using analysis of similarities (ANOSIM) and permutational multivariate analysis of variance (PERMANOVA) on Aitchison and/or Bray--Curtis distances~\cite{Li2019MiPro,braycurtis1957}. In the workflow, Aitchison distance is computed as Euclidean distance after centered log-ratio transformation, whereas Bray--Curtis dissimilarity is computed after conversion to relative abundance. ANOSIM reports a separation statistic $R$, where larger $R$ indicates stronger between-batch separation~\cite{Clarke1993ANOSIM}, whereas PERMANOVA reports an effect size $R^2$, corresponding to the fraction of distance variation explained by batch~\cite{anderson2001permanova}. The accompanying permutation $p$-values evaluate whether the observed batch association is larger than expected under random relabeling of the batch variable. We also generate dissimilarity heatmaps summarizing mean between-batch distances across batch pairs. More homogeneous values indicate weaker batch-driven separation.

\item \textbf{Variance-partitioning analyses.}
To quantify how compositional variation is apportioned among target (phenotype), batch, and their overlap, we used two multivariate variance-partitioning analyses under Aitchison geometry: partial redundancy analysis (pRDA) and principal variance component analysis (PVCA)~\cite{Wang2023PLSDABatch}. pRDA partitions variance into target, batch, intersection, and residual components; the intersection term reflects batch $\times$ target imbalance or confounding, with larger values indicating stronger confounding. PVCA yields an analogous Target/Batch/Intersection/Residual decomposition by applying a variance-components (mixed-effects) model to retained principal component (PC) scores and aggregating components across PCs using eigenvalue-based weights. We also computed per-feature one-way analysis of variance (ANOVA) $R^2$ values for batch and target and summarized them, for example by the median across features.

\item \textbf{Mixing and separation metrics based on neighborhood and embedding.}
We quantify local mixing using a $k$-nearest neighbor ($k$-NN) alignment score, where higher values indicate stronger cross-batch neighborhood mixing~\cite{Butler2018Seurat,Luecken2022scIB}. We compute a Uniform Manifold Approximation and Projection (UMAP) embedding~\cite{McInnes2018UMAP} and report an entropy-of-batch-mixing score on UMAP neighborhoods, where higher values indicate more uniform batch mixing~\cite{Haghverdi2018MNN,Luecken2022scIB}. To quantify phenotype preservation and separation, we report the silhouette coefficient with respect to target labels on the UMAP coordinates, where higher values indicate stronger separation~\cite{Rousseeuw1987Silhouette}.
\end{enumerate}

In the anaerobic digestion case study, these diagnostics were interpreted jointly rather than as a single ranking criterion. Ordination and distance-based tests summarize residual batch separation, variance-partitioning analyses describe how variation is attributed to batch, target and their overlap, and neighborhood metrics summarize local cross-batch mixing and target-label separation. Together, these outputs allow users to identify whether a correction output mainly reduces batch-associated structure, preserves target-associated structure, or changes both.


\subsection*{Supplementary Note 2. Input validation and data-cleaning boundary}
\label{sec:supp-validation}

mBatchNet is designed for processed microbiome abundance or count tables with matching sample metadata. It is not intended to replace upstream quality control, feature filtering, missing-value handling, or outlier review. During upload, the server validates core input requirements before preprocessing and writes a validation report for the session. Blocking errors include missing upload files, malformed numeric matrices, blank, non-numeric, NA, NaN or Inf matrix values, metadata row-count mismatch, missing selected batch or target metadata variables, insufficient batch or target levels, identical batch and target selections, all-zero sample rows, and matrices exceeding the public-server size limits. The server also reports warnings for large uploads, all-zero feature columns, high sparsity, negative or transformed-looking values, and strong batch-target association. mBatchNet does not impute missing values, remove outliers, or automatically decide which samples or features should be excluded; users should perform those data-cleaning steps before upload when they are needed.


\newpage
\section*{Supplementary Figures}
\label{sec:supp-figures}

\vspace*{\fill}
\begin{figure}[h!]
\begin{center}
\includegraphics[width=\textwidth]{FigSu1.pdf}
\end{center}
\textbf{Supplementary Figure S1. Sample composition by batch and target group.}
Proportion of samples in each batch stratified by the two target (phenotypic) cohorts. The anaerobic digestion example contains 75 samples and 231 OTUs. The two initial phenol concentration groups contain 26 and 49 samples, respectively, and the five processing-date batches contain 9, 16, 21, 17 and 12 samples. Batch-target composition is non-uniform across batches, which motivates inspecting batch-target balance before correction.
\end{figure}
\vspace*{\fill}

\newpage
\vspace*{\fill}
\begin{figure}[h!]
\begin{center}
\includegraphics[width=\textwidth]{FigSu2.pdf}
\end{center}
\textbf{Supplementary Figure S2. Aitchison-distance PCoA colored by target group.}
PCoA ordinations in Aitchison space, colored by the two target cohorts, shown before correction and after each batch-correction method.
\end{figure}
\vspace*{\fill}

\newpage
\vspace*{\fill}
\begin{figure}[h!]
\begin{center}
\includegraphics[width=\textwidth]{FigSu3.pdf}
\end{center}
\textbf{Supplementary Figure S3. PCA colored by batch.}
PCA embeddings colored by batch, shown before correction and after each batch-correction method.
\end{figure}
\vspace*{\fill}

\newpage
\vspace*{\fill}
\begin{figure}[h!]
\begin{center}
\includegraphics[width=\textwidth]{FigSu4.pdf}
\end{center}
\textbf{Supplementary Figure S4. PCA colored by target group.}
PCA embeddings colored by the two target cohorts, shown before correction and after each batch-correction method.
\end{figure}
\vspace*{\fill}

\newpage
\vspace*{\fill}
\begin{figure}[h!]
\begin{center}
\includegraphics[width=\textwidth]{FigSu5.pdf}
\end{center}
\textbf{Supplementary Figure S5. Aitchison-distance NMDS colored by batch.}
NMDS ordinations in Aitchison space colored by batch, shown before correction and after each batch-correction method.
\end{figure}
\vspace*{\fill}

\newpage
\vspace*{\fill}
\begin{figure}[h!]
\begin{center}
\includegraphics[width=\textwidth]{FigSu6.pdf}
\end{center}
\textbf{Supplementary Figure S6. Aitchison-distance NMDS colored by target group.}
NMDS ordinations in Aitchison space colored by the two target cohorts, shown before correction and after each batch-correction method.
\end{figure}
\vspace*{\fill}

\newpage
\vspace*{\fill}
\begin{figure}[h!]
\begin{center}
\includegraphics[width=\textwidth]{FigSu7.pdf}
\end{center}
\textbf{Supplementary Figure S7. PVCA variance partition across methods.}
Principal variance component analysis (PVCA) showing the fraction of variance attributable to target, batch, their overlap, and residual components for each method.
\end{figure}
\vspace*{\fill}

\newpage
\vspace*{\fill}
\begin{figure}[h!]
\begin{center}
\includegraphics[width=\textwidth]{FigSu8.pdf}
\end{center}
\textbf{Supplementary Figure S8. Batch-correction evaluation metrics across methods.}
Comparison of alignment, silhouette, and entropy-based metrics across methods, summarizing cross-batch mixing and target separation.
\end{figure}
\vspace*{\fill}


\newpage
\section*{Supplementary Tables}
\label{sec:supp-tables}

\begin{table}[H]
\centering
\caption{\textbf{Supplementary Table S1. Bray--Curtis batch-association statistics for the anaerobic digestion case study.} ANOSIM and PERMANOVA were computed using Bray--Curtis distances on relative-abundance profiles with 999 permutations.}
\label{tab:case-study-bray-stats}
\begin{tabular}{lrrrr}
\toprule
Method & ANOSIM $R$ & ANOSIM $p$ & PERMANOVA $R^2$ & PERMANOVA $p$ \\
\midrule
Before correction & 0.637 & 0.001 & 0.457 & 0.001 \\
ConQuR & -0.051 & 0.974 & 0.029 & 0.994 \\
DEBIAS-M & -0.052 & 0.981 & 0.035 & 0.908 \\
MetaDICT & 0.039 & 0.109 & 0.093 & 0.009 \\
MMUPHin & -0.010 & 0.648 & 0.069 & 0.148 \\
PLSDA-batch & 0.377 & 0.001 & 0.377 & 0.001 \\
RUV-III-NB & 0.639 & 0.001 & 0.460 & 0.001 \\
\bottomrule
\end{tabular}
\end{table}


\newpage
\section*{Supplementary References}
\label{sec:supp-refs}

\begin{thebibliography}{99}

\bibitem{Li2019MiPro}
Li, L., Abou-Samra, E., Ning, Z., Zhang, X., Mayne, J., Wang, J., Cheng, K., Walker, K., Stintzi, A. and Figeys, D.
An in vitro model maintaining taxon-specific functional activities of the gut microbiome.
\textit{Nature Communications}, 10:4146, 2019.

\bibitem{braycurtis1957}
Bray, J. R. and Curtis, J. T.
An ordination of the upland forest communities of southern Wisconsin.
\textit{Ecological Monographs}, 27(4):325--349, 1957.

\bibitem{Clarke1993ANOSIM}
Clarke, K. R.
Non-parametric multivariate analyses of changes in community structure.
\textit{Australian Journal of Ecology}, 18(1):117--143, 1993.

\bibitem{anderson2001permanova}
Anderson, M. J.
A new method for non-parametric multivariate analysis of variance.
\textit{Austral Ecology}, 26(1):32--46, 2001.

\bibitem{Wang2023PLSDABatch}
Wang, Y. and L{\^e} Cao, K.-A.
PLSDA-batch: a multivariate framework to correct for batch effects in microbiome data.
\textit{Briefings in Bioinformatics}, 24(2):bbac622, 2023.

\bibitem{Butler2018Seurat}
Butler, A., Hoffman, P., Smibert, P., Papalexi, E. and Satija, R.
Integrating single-cell transcriptomic data across different conditions, technologies, and species.
\textit{Nature Biotechnology}, 36(5):411--420, 2018.

\bibitem{Luecken2022scIB}
Luecken, M. D., B{\"u}ttner, M., Chaichoompu, K., Danese, A., Interlandi, M., Mueller, M. F., Strobl, D. C., Zappia, L., Dugas, M., Colom{\'e}-Tatch{\'e}, M. and Theis, F. J.
Benchmarking atlas-level data integration in single-cell genomics.
\textit{Nature Methods}, 19:41--50, 2022.

\bibitem{McInnes2018UMAP}
McInnes, L., Healy, J., Saul, N. and Gro{\ss}berger, L.
UMAP: Uniform Manifold Approximation and Projection.
\textit{Journal of Open Source Software}, 3(29):861, 2018.

\bibitem{Haghverdi2018MNN}
Haghverdi, L., Lun, A. T. L., Morgan, M. D. and Marioni, J. C.
Batch effects in single-cell RNA-sequencing data are corrected by matching mutual nearest neighbors.
\textit{Nature Biotechnology}, 36:421--427, 2018.

\bibitem{Rousseeuw1987Silhouette}
Rousseeuw, P. J.
Silhouettes: a graphical aid to the interpretation and validation of cluster analysis.
\textit{Journal of Computational and Applied Mathematics}, 20:53--65, 1987.

\end{thebibliography}

\end{document}
